\documentclass[11pt]{article}
\usepackage[a4paper,margin=2cm]{geometry}
\usepackage[ngerman]{babel}\usepackage[T1]{fontenc}\usepackage[utf8]{inputenc}
\usepackage{amsmath,amssymb,tikz,xcolor}\pagestyle{empty}\setlength\parindent{0pt}
\usetikzlibrary{calc,angles}
\newcommand{\karo}[1]{\par\smallskip\begin{tikzpicture}\draw[step=5mm,gray!30,very thin] (0,0) grid (17,#1);\end{tikzpicture}\par}
\newcommand{\nr}[1]{\par\bigskip\textbf{#1}\quad}
% \tri{scale}{R}{P}{Q}{lab RP}{lab RQ}{lab PQ}  rechter Winkel bei R
\newcommand{\tri}[7]{\begin{tikzpicture}[scale=#1,thick]
\coordinate (R) at (#2);\coordinate (P) at (#3);\coordinate (Q) at (#4);
\coordinate (G) at ($1/3*(R)+1/3*(P)+1/3*(Q)$);
\draw (R)--(P)--(Q)--cycle;
\pic[draw,thin,angle radius=2.2mm]{right angle=P--R--Q};
\node at ($(R)!0.5!(P)!-0.7!(G)$) {\small #5};
\node at ($(R)!0.5!(Q)!-0.7!(G)$) {\small #6};
\node at ($(P)!0.5!(Q)!-0.75!(G)$) {\small #7};
\end{tikzpicture}}
\begin{document}
{\Large\bfseries Hypotenuse berechnen}\quad{\color{gray}Klasse 9}\par\medskip

\nr{1.}Finde die Hypotenuse. Kreise ihren Buchstaben ein.\par\medskip
\begin{tabular}{cccc}
\tri{0.5}{0,0}{5,0}{0,3.5}{$p$}{$q$}{$r$} &
\tri{0.5}{5,0}{0,0}{5,4}{$z$}{$x$}{$y$} &
\tri{0.5}{4,3}{0,2}{5,0.6}{$k$}{$m$}{$l$} &
\begin{tikzpicture}[scale=0.5,thick]\draw (0,0)--(5,0)--(1.6,3.2)--cycle;
\node at (2.5,-0.5){\small $d$};\node at (3.7,1.9){\small $e$};\node at (0.4,1.8){\small $f$};\end{tikzpicture}\\
a) & b) & c) & d)
\end{tabular}

\nr{2.}Notiere den Satz des Pythagoras. Verwende eine Skizze. Schlage im Tafelwerk nach und trage die Seite ein: S.~\rule{1.2cm}{0.4pt}
\karo{3}

\nr{3.}Stelle den Satz des Pythagoras auf. Rechne nicht.\par\medskip
\begin{tabular}{ccc}
\tri{0.45}{0,0}{5,0}{0,3}{$u$}{$v$}{$w$} &
\tri{0.45}{3,0}{0,0.8}{3.9,4}{$t$}{$s$}{$r$} &
\tri{0.45}{4,4}{0,4}{4,0}{$g$}{$e$}{$f$}\\[1mm]
a) \rule{3.5cm}{0.4pt} & b) \rule{3.5cm}{0.4pt} & c) \rule{3.5cm}{0.4pt}
\end{tabular}

\nr{4.}Berechne die fehlende Seite.\par\medskip
\begin{tabular}{cccc}
\tri{0.5}{0,0}{4,0}{0,3}{8 cm}{6 cm}{$c$} &
\tri{0.45}{0,0}{6,0}{0,2.5}{12 cm}{5 cm}{$c$} &
\tri{0.45}{4.5,2.4}{0,2.4}{4.5,0}{15 cm}{8 cm}{$x$} &
\tri{0.5}{0,0}{0,4}{3.2,0}{24 cm}{$a$}{25 cm}\\
a) & b) & c) & d)
\end{tabular}\par\smallskip
{\color{gray}a)\quad$\begin{aligned}[t]
c^2 &= 6^2 + 8^2\\
c^2 &= 100 &&\big|\ \sqrt{\ }\\
c &= 10\ \text{cm}
\end{aligned}$}
\karo{5}

\newpage
\nr{5.}Berechne die Hypotenuse. Runde auf eine Stelle nach dem Komma.\par\medskip
\begin{tabular}{ccc}
\tri{0.5}{0,0}{4,0}{0,3}{12 cm}{9 cm}{$c$} &
\tri{0.5}{3.5,0}{3.5,2.4}{0,0}{2,4 cm}{3,5 cm}{$c$} &
\tri{0.5}{0,0}{4.8,0}{0,2}{1,2 m}{50 cm}{$c$}\\
a) & b) & c)
\end{tabular}
\karo{5}

\nr{6.}Eine Leiter lehnt an einer Wand. Ihr Fuß steht 0,9 m von der Wand weg, oben reicht sie 1,2 m hoch. Wie lang ist die Leiter?\par\medskip
\begin{tikzpicture}[scale=1.3,thick]
\draw (0,0)--(0,1.6); \draw (-0.3,0)--(1.4,0); \draw (0,1.2)--(0.9,0);
\draw[thin] (0,0.18)--(0.18,0.18)--(0.18,0);
\node[left] at (0,0.6){\small 1,2 m};\node[below] at (0.45,0){\small 0,9 m};
\end{tikzpicture}
\karo{2.5}
\nr{7.}Ein rechteckiger Park ist 160 m lang und 70 m breit. Lena will von einer Ecke zur gegenüberliegenden Ecke. Bisher ging sie außen am Rand entlang. Jetzt gibt es einen geraden Weg quer durch den Park. Wie viele Meter spart sie?\par\medskip
\begin{tikzpicture}[scale=0.04,thick]
\draw[fill=green!8] (0,0) rectangle (160,70); \draw[dashed] (0,0)--(160,70);
\node[below left] at (0,0){\small Start};\node[above right] at (160,70){\small Ziel};
\node[below] at (80,0){\small 160 m};\node[right] at (160,35){\small 70 m};
\end{tikzpicture}
\karo{3.5}

\nr{8.}Ein Schrank ist 240 cm hoch und 70 cm tief. Er wird liegend zusammengebaut und dann aufgerichtet; dabei kippt er über eine untere Kante. Das Zimmer ist 245 cm hoch. Geht das? Begründe mit einer Rechnung.\par\medskip
\begin{tikzpicture}[scale=0.012,thick]
\draw (0,0)--(420,0); \draw[dashed] (0,245)--(420,245); \node[above] at (330,245){\small Decke: 245 cm};
\draw (40,0) rectangle (110,240); \node[below] at (75,0){\small 70 cm}; \node[left] at (40,120){\small 240 cm};
\draw[->,thin] (200,80) arc (0:60:60); \node at (300,100){\small aufrichten};
\end{tikzpicture}
\karo{4}
\vfill
{\color{gray}\footnotesize Lösungen: 4: a) 10 cm; b) 13 cm; c) 17 cm; d) $a$ ist Kathete: $a^2 = 25^2 - 24^2 = 49$, $a = 7$ cm \quad 5: a) 15 cm; b) 4,2 cm; c) 1,3 m \quad 6: 1,5 m \quad 7: etwa 55,4 m \quad 8: Nein, Diagonale 250 cm $>$ 245 cm}
\end{document}
